\chapter{Obsessive Compulsive Disorder (OCD) Test}
\section*{What Is This Test?} OCD assessment identifies intrusive thoughts, urges, or images and repetitive behaviors or mental acts that cause distress or impairment.\section*{Alternative Names} OCD screening; obsessive-compulsive assessment; Y-BOCS assessment.\section*{Principle} Interview and validated scales assess obsessions, compulsions, time burden, distress, insight, and function.\section*{Why This Test Is Ordered} It is used for recurrent unwanted thoughts, rituals, avoidance, or anxiety interfering with daily life.\section*{Primary vs Secondary Test} Clinical interview is primary; questionnaires support but do not diagnose alone.\section*{How This Test Is Performed} The patient describes symptoms, onset, triggers, avoidance, comorbidities, medicines, and safety concerns.\section*{What Preparation Is Needed} None; bring symptom history and prior treatment information.\section*{Understanding Results} Symptoms must be distinguished from generalized anxiety, depression, psychosis, autism, tic disorders, or substance effects.\section*{Follow-up Tests} Psychiatric evaluation, medical review, substance or sleep assessment, and structured severity scales may follow.\section*{References} \begin{enumerate}\item MedlinePlus. Obsessive Compulsive Disorder (OCD) Test.\item International OCD Foundation. Assessment resources.\end{enumerate}\clearpage\section*{Understanding Results and Follow-up}
The laboratory report should identify the specimen, method, units, reference interval or decision limit, and whether the result is positive, negative, abnormal, or inconclusive. Reference intervals vary with age, sex, pregnancy, specimen type, assay, and laboratory; a result outside a range is not automatically a diagnosis, and a result within a range does not always exclude disease. Discuss unexpected results with the ordering clinician, who can decide whether repeat or confirmatory testing, treatment, or referral is appropriate. Do not change medicines or delay urgent care based on an isolated result.
\section*{Regulatory Status and Authorization Date}This chapter describes a test category, not a specific commercial product. FDA, CDSCO, EU IVDR, and MHRA approval or registration dates therefore cannot be assigned without the assay manufacturer, product name, instrument, and intended use. For laboratory-developed tests, an individual agency approval date may not exist. See the Preface for the interpretation of regulatory status.
\clearpage
